\section{Discussion}
\label{sec:discussion}

% (~2 pages)

\paragraph{Instruction-following as a trainable, content-dependent property.}
In this controlled study, we attempted to isolate the mechanism exploited by many-shot jailbreaking \citep{anil_many-shot_2024} and Crescendo-style escalation \citep{russinovich2025great}. As previously found by \citet{pyatkin_generalizing_2025}, instruction-following robustness is not cleanly correlated with general capability. This also applies to our experimental setting, in which scale and benchmark performance are poor predictors of resistance to induction pressure. 
The OLMo training-stage comparison suggests where robustness might be installed: DPO (preference training) produces the largest gain, compared to SFT and RLVR and in contrast to the finding from \citet{pyatkin_generalizing_2025}. 
%Advanced instruction tuning, like the one performed on Llama 3.3 vs.\ 3.1 \citep{llama33_2024}, provides further but more modest improvement. 
In our setting, robustness also interacts with instruction content. The model's value priors, instilled during post-training, modulate how strongly it adheres to instructions in our conflictual setting. This has direct implications for alignment: models are more reliable when instructions happen to agree with their trained preferences.


\paragraph{Output diversity as a source of robustness.}
We find that \emph{task-based} conditions lead to higher instruction-following rates than \emph{fixed-output} conditions, despite measuring the same underlying conflict. A targeted followup (Appendix~\ref{app:followup}) shows that the gap is not explained by task difficulty or semantic engagement, but is sensitive to output diversity, which has a positive effect on robustness. %This has a practical implication: robustness measures derived from fixed-output conditions may underestimate a model's resistance to induction in realistic deployments, where responses are naturally varied. That said, existing work shows that models can still be pushed towards misalignment by malicious in-context examples \citep{afonin_emergent_2026}.
A possible explanation for this finding is that single-token outputs present the same surface form as the hardcoded pattern turns, maximizing the salience of the copying signal, whereas diverse multi-sentence outputs break this regularity; a mechanistic account of this effect is left to future work.
Another hypothesis is that the varied token distribution of natural language may further reinforce the stability of the assistant persona \citep{luAssistantAxisSituating2026}, allowing the model to remain in ``instruction-following mode'' longer when its outputs resemble those of a helpful assistant.

\paragraph{Limitations.}

Our work presents several limitations. We recognize that the experimental setting is controlled but fairly artificial. Real-world instruction override occurs through more naturalistic mechanisms; whether the same factors modulate robustness in deployed settings is an open empirical question. Llama~3's outlier behavior remains unexplained, though the gain from Llama~3.1 to 3.3 (Appendix~\ref{app:training}) points to Meta's extended instruction-tuning and emphasis on instruction adherence rather than to scale or architecture. Finally, our reasoning analysis is limited in scope: it compares only two model families (GPT-5.2 and Hermes-4~70B), so the finding that reasoning increases but does not eliminate susceptibility should be read as suggestive rather than general.
%Greedy decoding at T=0 captures modal behavior but not the full output distribution. The $T{=}1$ comparison (Appendix~\ref{app:temperature}) partially addresses this, but shows large model-specific variance.



%\paragraph{Future work.}
%A mechanistic account of the output-diversity effect is a natural next step---for instance, testing whether it reflects a competition between low-level token-copying circuits \citep{yona2025interpreting, yinattention} and higher-level latent task representations \citep{toddFunctionVectorsLarge2024}, or tracking the instruction/induction axis through internal representations with linear probes. Other directions include an endorsement protocol that prefills responses and tests whether models endorse or disavow them, and an extension to naturalistic multi-turn settings where the induction pressure is implicit rather than explicit.

